% ============================================================================
% CONCLUSION G\'eN\'eRALE
% source : intro_conclu_gen/conclusion_generale.tex
% Version : 2.2 — alignée sur la formalisation Lean4 révisée
% ============================================================================

%\chapter*{CONCLUSION G\'eN\'eRALE}
%\addcontentsline{toc}{chapter}{Conclusion g\'en\'erale}
\label{chap:conclusion_generale}

% ----------------------------------------------------------------------------
\section{RETOUR SUR LES OBJECTIFS DE RECHERCHE}
\label{sec:cg_questions}
% ----------------------------------------------------------------------------
\begin{sloppypar}

Le \textbf{paradoxe de l'authenticit\'e invisible}~\cite{MINKA2026CLO}, selon lequel une preuve
cryptographique math\'ematiquement irr\'eprochable peut demeurer institutionnellement
opaque, formalise une tension latente entre les exigences cryptographiques et
l'utilit\'e institutionnelle qui passe par l'interpr\'etabilit\'e des preuves par
les arbitres (juge, r\'egulateur, auditeur) non-experts. Nous nous sommes servi
d'exemples illustratifs~\cite{IACR2025incident,Adida09UCL} d\'ecrits \`a la
section~\ref{subsec:cas_etude} pour mat\'erialiser la r\'ealit\'e de cette
situation et en faciliter la compr\'ehension. Ce paradoxe g\'en\`ere la
probl\'ematique centrale formul\'ee ci-dessous.

\begin{boiteProblematique}{Probl\'ematique centrale}
\textit{Comment, d'une part, concilier la s\'ecurit\'e cryptographique qui
requiert l'opacit\'e computationnelle avec l'interpr\'etabilit\'e
institutionnelle qui requiert la transparence s\'emantique~? D'autre part,
existe-t-il des limites fondamentales aux compromis r\'ealisables entre
confidentialit\'e, fiabilit\'e et opposabilit\'e~?}
\end{boiteProblematique}

Cette th\`ese y r\'epond en traitant les six objectifs qui en
d\'ecoulent, avec le statut propre \`a chacun.

\medskip

\textbf{Obj~\ref{obj:obj1_unification}. Formalisation.} Peut-on formaliser
l'utilit\'e institutionnelle d'un protocole cryptographique~? Nous avons
d\'efini trois mesures issues de la th\'eorie de
l'information~\cite{Shannon48,CoverThomas06}, \`a savoir $\Conf$ (statistique), la
fiabilit\'e distingu\'ee en $\Relrec$ (reconstruction) et $\Relver$
(v\'erification), enfin $\Opp = I(V;\Sigma|C)/H(V|C)$, dont l'in\'egalit\'e
$I(V;\Sigma|C) \leq \min\{H(\Sigma|C),H(V|C)\} \leq \alpha H(\mathcal{C})$ formalise
le goulot contextuel.
L'opposabilit\'e $\Opp$ \'etait absente de tout cadre cryptographique
ant\'erieur.

\textbf{Obj~\ref{obj:obj2_impossibilite}. Bornes.} Existent-ils des limites
fondamentales sur les compromis r\'ealisables~? Sous quatre hypoth\`eses
institutionnelles explicites, nous avons \'etabli la borne
\begin{equation}
  \Conf \cdot \Relrec \cdot \Opp \;\leq\;
  \frac{1}{\gamma}\left(\frac{\Delta}{\lambda} + \frac{\kappa(\Delta+1)}{\lambda^2} + \frac{\kappa^2}{\lambda^3}\right)
\end{equation}
o\`u $\gamma = \alpha H(C) - 1 \geq 1$ et $\kappa$ est la constante du
r\'egime de pr\'eservation de la confidentialit\'e
(Hypoth\`ese~\ref{ass:regime}). Le terme dominant de cette borne est
$\Delta/(\gamma\lambda)$~; les autres sont en $O(\lambda^{-2})$. La forme raffin\'ee
\eqref{eq:trilemme_cro_raffine} identifie le r\'egime ($H(\mathcal{C}) > H^\sharp$) o\`u la
capacit\'e d'extraction institutionnelle est la contrainte liante, la puissance
de calcul de l'adversaire passant au second plan. La formulation
publi\'ee~\cite{MINKA2026CRO} s'obtient comme cas particulier.

\textbf{Obj~\ref{obj:obj3_postquantique}. Extension quantique.} Les bornes CRO
se pr\'eservent-elles face aux adversaires quantiques~? Nous avons \'etabli
le Trilemme CRO post-quantique
\begin{equation}
  \Conf^q \cdot \Relrec^q \cdot \Opp^q \;\leq\;
  \frac{1}{\gamma_q}\left(\frac{\Delta_q}{\lambda} + \frac{\kappa(\Delta_q+1)}{\lambda^2} + \frac{\kappa^2}{\lambda^3}\right)
\end{equation}
o\`u $\Delta_q$ est la capacit\'e d'extraction quantique et $\eta_q$ quantifie
la perte d'interpr\'etabilit\'e, nulle pour les preuves qui sont des cha\^{\i}nes
de bits classiques, dont les signatures ML-DSA et SLH-DSA.

\textbf{Obj~\ref{obj:obj4_saturation}. Optimalit\'e.} Peut-on approcher ces
bornes~? La famille SOTP propose une construction par r\'egime, SOTP-Low,
SOTP-Base (projection Johnson-Lindenstrauss) et SOTP-High (Groth16). Leur
confidentialit\'e et leur fiabilit\'e de v\'erification sont prouv\'ees sous
des hypoth\`eses standard~; leur opposabilit\'e est guid\'ee par une heuristique
de conception, tandis que leur optimalit\'e demeure conjecturale
(Conjecture~\ref{conj:optimalite_universelle}). Cet objectif est donc
partiellement atteint.

\textbf{Obj~\ref{obj:obj5_verification}. V\'erification.} Comment certifier la
correction d'un th\'eor\`eme complexe~? La formalisation Lean4 a
certifi\'e m\'ecaniquement le noyau alg\'ebrique du th\'eor\`eme
(section~\ref{sec:lean4}), ainsi que la forme raffin\'ee de la borne. Fano et
Pinsker y sont pos\'es comme propri\'et\'es des instances de protocole~; les
hypoth\`eses institutionnelles y sont aussi explicit\'ees. La formalisation a produit
deux t\'emoins, l'un de coh\'erence (un point du mod\`ele abstrait o\`u la
borne est atteinte, non r\'ealisable par un protocole), l'autre de port\'ee
(une instance hors r\'egime atteint $1/4$). Le d\'etail de ce que la certification a rendu explicite, hypoth\`eses,
constantes et r\'egime de validit\'e, est pr\'esent\'e \`a la
sous-section~\ref{subsec:lean4_erreurs}.

\textbf{Obj~\ref{obj:obj6_experimentation}. Exp\'erimentation.} Comment rendre
ces bornes testables~? Six protocoles op\'erationnels dotent chacun des
param\`etres institutionnels d'une d\'efinition mesurable, d'un estimateur et
d'un plan d'\'echantillonnage. Ils fixent aussi les seuils \`a partir desquels le
mod\`ele serait r\'efut\'e (tableau~\ref{tab:criteres_refutation}). Une
calibration illustrative sur des syst\`emes d\'eploy\'es v\'erifie que les ordres de
grandeur retenus sont mutuellement coh\'erents~; ces syst\`emes ne sont pas des
instances du th\'eor\`eme. Le programme exp\'erimental lui-m\^eme n'a d'ailleurs pas \'et\'e
d\'eploy\'e dans le cadre de cette th\`ese.

\medskip
\`A la question de savoir si confidentialit\'e, fiabilit\'e et opposabilit\'e
peuvent \^etre simultan\'ement maximis\'ees, la th\`ese r\'epond par la n\'egative
sous quatre hypoth\`eses institutionnelles explicites. Le produit y est born\'e
par une quantit\'e dont le terme dominant est $\Delta/(\gamma\lambda)$, tandis que la
forme raffin\'ee identifie exactement le r\'egime ($H(\mathcal{C}) > H^\sharp$) o\`u la
capacit\'e d'extraction institutionnelle est la contrainte liante. Le noyau
alg\'ebrique de ce r\'esultat est certifi\'e en Lean4 et sa port\'ee est
d\'elimit\'ee, puisque les syst\`emes d\'eploy\'es \`a haute fiabilit\'e de
v\'erification n'en sont pas des instances et servent \`a calibrer les
ordres de grandeur. L'existence de constructions atteignant la borne est
formul\'ee comme conjecture et constitue, avec la mesure exp\'erimentale de
$\Delta$ par l'estimateur $\hat{I}(V;\Sigma|M,C)$, le programme ouvert par ce travail.

\end{sloppypar}

% ----------------------------------------------------------------------------
\section{PRINCIPALES CONTRIBUTIONS}
\label{sec:cg_contributions}
% ----------------------------------------------------------------------------
\begin{sloppypar}

Le cycle Obj~\ref{obj:obj1_unification}~$\to$~Obj~\ref{obj:obj2_impossibilite}~$\to$~Obj~\ref{obj:obj3_postquantique}~$\to$~Obj~\ref{obj:obj4_saturation}~$\to$~Obj~\ref{obj:obj5_verification}~$\to$~Obj~\ref{obj:obj6_experimentation} suivant la progression
F$\to$B$\to$E$\to$S$\to$C$\to$X
constitue le squelette logique de la th\`ese. Six contributions
techniques (C1--C6) et trois apports transversaux en \'emergent.

\subsection*{Contributions techniques (C1--C6)}

\begin{table}[htbp]
\centering
\caption{Synth\`ese des six contributions}
\label{tab:synthese_contributions}
\small
\begin{tabular}{@{}llp{6cm}@{}}
\toprule
\textbf{\'Etape} & \textbf{Contribution} & \textbf{R\'esultat cl\'e} \\
\midrule
F (Formaliser) & C1~: Cadre CLO & M\'etriques $\Conf$, $\Relrec$/$\Relver$, $\Opp$~; goulot contextuel \\
B (Borner) & C2~: Trilemme CRO classique & Borne $\gamma^{-1}(\Delta/\lambda + \kappa(\Delta+1)/\lambda^2 + \kappa^2/\lambda^3)$ et forme raffin\'ee \\
E (\'Etendre) & C3~: Trilemme CRO post-quantique & Borne avec $\Delta_q$, $\eta_q$~; calibration NIST PQC \\
S (Saturer) & C4~: Famille SOTP & Trois constructions ($\Conf$, $\Relver$ prouv\'ees)~; optimalit\'e conjectur\'ee \\
C (Certifier) & C5~: Certification Lean4 & Noyau v\'erifi\'e, t\'emoins de coh\'erence et port\'ee \\
X (eXp\'erimenter) & C6~: Programme de validation & Protocoles exp\'erimentaux sp\'ecifi\'es, falsifiables \\
\bottomrule
\end{tabular}
\end{table}

\textbf{C1. Cadre CLO}~\cite{MINKA2026CLO}. Premier cadre unifi\'e
formalisant $\Conf$, $\Rel$ et $\Opp$ via la th\'eorie de l'information, avec
la distinction entre fiabilit\'e de reconstruction et de v\'erification.
Le th\'eor\`eme du goulot d'\'etranglement contextuel \'etablit que
$I(V;\Sigma|C) \leq \alpha H(\mathcal{C})$~; un verdict d\'eterministe donne $\Opp = 1$.

\textbf{C2. Trilemme CRO conditionnel}~\cite{MINKA2026CRO}. Premi\`ere
borne prouv\'ee sur les compromis $(\Conf, \Rel, \Opp)$, conditionnelle
aux param\`etres $(\Delta, \alpha, \beta)$ et au r\'egime de
confidentialit\'e $I(M;\Sigma|C) \leq \kappa/\lambda$. La forme explicite
est $\gamma^{-1}(\Delta/\lambda + \kappa(\Delta+1)/\lambda^2 + \kappa^2/\lambda^3)$~;
la forme raffin\'ee, fond\'ee sur l'identit\'e
$I(V;\Sigma|C) = I(V;M|C) + I(V;\Sigma|M,C)$ pour un interpr\`ete admissible, distingue
un r\'egime o\`u Fano limite l'opposabilit\'e et un r\'egime o\`u $\Delta$ la limite.
La finesse de la borne reste une question ouverte.

\textbf{C3. Trilemme CRO post-quantique}~\cite{MINKA2026PQCRO}. Extension
aux adversaires quantiques avec param\`etres $\Delta_q$ et $\eta_q$.
Calibration illustrative sur ML-DSA et SLH-DSA, pour lesquelles $\eta_q = 0$.

\textbf{C4. Famille SOTP}. Trois constructions, une par r\'egime~;
confidentialit\'e et fiabilit\'e de v\'erification prouv\'ees sous PRG, DLog et
q-PKE~; opposabilit\'e guid\'ee par une heuristique de conception~; optimalit\'e
conjectur\'ee. Introduction du paradigme Gold-Iron-Clay.

\textbf{C5. Formalisation Lean4}. Certification m\'ecanis\'ee
du noyau alg\'ebrique et de la forme raffin\'ee. Fano et Pinsker pos\'es comme
propri\'et\'es d'instance, explicitation des hypoth\`eses institutionnelles
et production de deux t\'emoins, l'un de coh\'erence (mod\`ele abstrait), l'autre de
port\'ee. Les formulations alternatives issues de
la certification sont pr\'esent\'ees \`a la sous-section~\ref{subsec:lean4_erreurs}.

\textbf{C6. Programme de validation falsifiable}. Les protocoles
exp\'erimentaux permettant de mesurer $\Delta$, $\alpha$, $\beta$,
$\Delta_q$ et $\eta_q$ sont sp\'ecifi\'es jusqu'au niveau
op\'erationnel~(chapitre~\ref{chap:experimentation}).

\subsection*{Apports transversaux}

\textbf{Apport constructif (de la borne au guide).} La famille SOTP
transforme le Trilemme CRO en guide de conception, sans \'etablir que la
borne est atteinte, question qui fait l'objet de la
Conjecture~\ref{conj:optimalite_universelle}. Comme CAP trouve son sens dans
les syst\`emes distribu\'es r\'eels et FLP dans les protocoles partiels, le
Trilemme CRO trouve un prolongement dans les trois variantes
SOTP~(section~\ref{subsec:portee_constructive}).

\textbf{Apport m\'ethodologique (la le\c{c}on de la formalisation).}
La formalisation Lean4 a conduit \`a des formulations plus pr\'ecises de trois
\'etapes des preuves publi\'ees (sens d'une in\'egalit\'e, statut de
l'Hypoth\`ese~2, r\'esolution d'une optimisation). Elle a \'etabli par un cas
limite la n\'ecessit\'e du r\'egime de confidentialit\'e. La le\c{c}on, selon
laquelle la formalisation gagne \`a pr\'ec\'eder la publication, est d\'evelopp\'ee
\`a la sous-section~\ref{subsec:verification_corrige_lecon}. Les formulations
alternatives sont pr\'esent\'ees \`a la
sous-section~\ref{subsec:lean4_erreurs}.

\textbf{Apport \'epist\'emique (un programme falsifiable).} Le
chapitre~\ref{chap:experimentation} sp\'ecifie comment tester le
cadre, transformant un r\'esultat conditionnel en programme scientifique
falsifiable au sens de Popper~\cite{Popper1934}. Sans cette
sp\'ecification, le Trilemme CRO serait non testable~; avec elle, il est
testable et critiquable~(section~\ref{subsec:statut_protocoles_experimentaux}).

\end{sloppypar}

% ----------------------------------------------------------------------------
\section{POSITIONNEMENT DANS LA LITT\'ERATURE}
\label{sec:cg_positionnement}
% ----------------------------------------------------------------------------
\begin{sloppypar}

Quatre corpus ont abord\'e des aspects partiels du probl\`eme, \`a savoir la
QIF~\cite{Smith09}, la s\'ecurit\'e utilisable~\cite{Bonneau12}, les \'etudes
juridiques~\cite{Kerr04} et les trilemmes d'impossibilit\'e
(CAP~\cite{Brewer00,GilbertLynch02}, FLP~\cite{FLP85},
Blockchain~\cite{Nakamoto08,Buterin17}). Le cadre CLO-CRO s'en distingue
en ce qu'il formalise simultan\'ement les trois dimensions
$(\Conf, \Rel, \Opp)$, prouve une borne quantitative conditionnelle
classique et quantique, propose des constructions dont l'optimalit\'e est
conjectur\'ee et certifie m\'ecaniquement une partie du r\'esultat.

\begin{table}[htbp]
\centering
\caption{Positionnement par rapport aux travaux existants.}
\label{tab:positionnement_travaux}
\resizebox{\textwidth}{!}{%
\begin{tabular}{@{}lllllll@{}}
\toprule
\textbf{Cadre} & \textbf{Conf} & \textbf{Rel} & \textbf{Opp} & \textbf{Borne} & \textbf{Optimalit\'e} & \textbf{Validation} \\
\midrule
QIF                    & Total     & Total     & Absent      & Partielle & --- & --- \\
S\'ecurit\'e utilisable    & Partielle & Absent    & Qualitative & Absent    & --- & Empirique \\
\'Etudes juridiques      & Absent    & Absent    & Qualitative & Absent    & --- & --- \\
CAP/FLP/Blockchain     & Partielle & Partielle & Absent      & Total     & Oui & --- \\
\textbf{CLO-CRO (ce travail)} & \textbf{Total} & \textbf{Total} & \textbf{Total} & \textbf{Oui (conditionnel)} & \textbf{Partiel (conjecture)} & \textbf{Sp\'ecifi\'e, non d\'eploy\'e} \\
\bottomrule
\end{tabular}}

\vspace{0.4em}
\footnotesize
\textit{\'Echelle de couverture.} Total (traitement formel complet)~;
Partielle (traitement formel restreint \`a certains cas)~; Qualitative
(traitement non formalis\'e)~; Absent (dimension non trait\'ee).
\end{table}

\end{sloppypar}

% ----------------------------------------------------------------------------
\section{PORT\'EE DE LA TH\`ESE}
\label{sec:cg_portee}
% ----------------------------------------------------------------------------
\begin{sloppypar}

La port\'ee de cette th\`ese est quadruple.

\textbf{Th\'eorique.} Les bornes C1--C3 \'etablissent des limites
fondamentales conditionnelles \`a des param\`etres institutionnels
mesurables.

\textbf{Constructive.} La famille SOTP (C4) propose, pour chaque r\'egime, une
construction aux propri\'et\'es cryptographiques prouv\'ees~; qu'elle approche
les bornes \`a facteur constant est conjectur\'e.

\textbf{Certifi\'ee.} La formalisation Lean4 (C5) certifie
m\'ecaniquement le noyau alg\'ebrique du Trilemme CRO classique.

\textbf{M\'ethodologique.} Le programme de validation (C6) et la
le\c{c}on de la formalisation
(section~\ref{subsec:verification_corrige_lecon}) fournissent \`a la
communaut\'e des outils et des m\'ethodes transf\'erables.

\end{sloppypar}

% ----------------------------------------------------------------------------
\section{LIMITATIONS ET TRAVAUX FUTURS}
\label{sec:cg_limitations}
% ----------------------------------------------------------------------------
\begin{sloppypar}

Quatre limitations d\'elimitent le p\'erim\`etre des contributions. Les
hypoth\`eses institutionnelles $(\Delta, \alpha, \beta)$ ne sont pas
encore valid\'ees par des \'etudes exp\'erimentales~; les valeurs retenues,
de quelques dizaines de bits, sont une calibration illustrative rendue
plausible par la litt\'erature sur la rationalit\'e born\'ee. L'opposabilit\'e
des constructions SOTP est heuristique et leur optimalit\'e conjectur\'ee. La formalisation Lean4 couvre le noyau
alg\'ebrique~; les constructions SOTP et l'extension post-quantique restent
hors de son p\'erim\`etre. Les constructions SOTP disposent d'impl\'ementations
de r\'ef\'erence dans le d\'ep\^ot compagnon, sans d\'eploiement en conditions
r\'eelles. Une
analyse compl\`ete des limitations est pr\'esent\'ee \`a la
section~\ref{subsec:limitations_problemes_ouverts}.

Les travaux futurs par contribution et par horizon sont d\'etaill\'es au
tableau~\ref{tab:perspectives_horizons}.

Le programme futur s'organise selon trois directions structurantes.

\textbf{Validation empirique.} Les param\`etres $\Delta$, $\alpha$,
$\beta$, $\Delta_q$, $\eta_q$ doivent \^etre mesur\'es sur des acteurs
institutionnels r\'eels. Le programme du
chapitre~\ref{chap:experimentation} fournit la feuille de route.

\textbf{Extension de la formalisation.} La formalisation Lean4 doit
\^etre \'etendue \`a SOTP et au cadre post-quantique, une fois les
biblioth\`eques de th\'eorie de l'information quantique disponibles dans
Mathlib~(section~\ref{subsubsec:extension_lean4_sotp_pq}).

\textbf{Extension th\'eorique.} Les variantes du Trilemme CRO pour le
chiffrement homomorphe, le calcul multi-parties et les contextes
dynamiques restent \`a d\'evelopper~(section~\ref{subsec:axe3_extensions}).

\end{sloppypar}

% ----------------------------------------------------------------------------
\section{SIGNIFICATION}
\label{sec:cg_signification}
% ----------------------------------------------------------------------------
\begin{sloppypar}

La cryptographie classique s'attache \`a la probabilit\'e qu'un
adversaire casse un protocole, alors que la cryptographie institutionnelle
s'interroge sur la probabilit\'e qu'un juge comprenne une
preuve. Loin d'\^etre une m\'etrique de confort
utilisateur, l'opposabilit\'e $\Opp$ traduit une contrainte propre \`a la nature de l'information,
de m\^eme rang que la confidentialit\'e. Le Trilemme CRO dit que,
sous ses hypoth\`eses, certains compromis sont impossibles, ind\'ependamment
du talent des ing\'enieurs.

La certification Lean4 illustre un pari \'epist\'emologique selon lequel les
th\'eor\`emes importants m\'eritent une v\'erification m\'ecanique.

La cryptographie institutionnelle dispose d\'esormais d'un cadre
compos\'e de m\'etriques (CLO), d'une borne d'impossibilit\'e conditionnelle classique et
quantique (CRO), de constructions (SOTP), d'une v\'erification formelle (Lean4)
et d'un programme de validation (C6). Le travail restant (finesse de la borne,
optimalit\'e des constructions, calibration des param\`etres,
extensions FHE/MPC, validation empirique, standardisation) constitue un
programme de recherche ouvert.

La cryptographie n'existe pas dans le vide. Elle sert
des institutions humaines soumises \`a des contraintes cognitives,
organisationnelles et proc\'edurales. L'ignorance de ces
contraintes engendre le paradoxe de l'authenticit\'e invisible.

\end{sloppypar}

% ----------------------------------------------------------------------------
\section*{MOT DE LA FIN}
\label{sec:cg_mot_fin}
% ----------------------------------------------------------------------------
\begin{sloppypar}

Cette th\`ese a transform\'e en cadre quantitatif l'intuition qualitative
selon laquelle les preuves cryptographiques sont opaques pour les interpr\`etes
institutionnels. Ce cadre r\'eunit des m\'etriques d\'efinies, des bornes
prouv\'ees sous hypoth\`eses, des constructions, une certification m\'ecanique
et un programme de validation falsifiable. Dans le r\'egime que la forme
raffin\'ee identifie, la contrainte dominante devient la capacit\'e d'extraction de l'interpr\`ete, la
puissance de calcul de l'adversaire passant au second plan.
Cette inversion de perspective est l'id\'ee centrale de cette th\`ese~; sa
confirmation empirique reste \`a apporter.

L'arbitre le rappelait, \emph{«~La rigueur ne se suppose pas. Elle se
v\'erifie.~»} Cette th\`ese a v\'erifi\'e et pr\'ecis\'e ses r\'esultats, document\'e
ces pr\'ecisions et sp\'ecifi\'e comment tester ce qui reste \`a
\'etablir. En cela, la science se distingue de la sp\'eculation.

\end{sloppypar}

% ----------------------------------------------------------------------------
% Bo\^ite finale
% ----------------------------------------------------------------------------
\begin{boiteArbitre}
L'arbitre referme le dossier sans avoir r\'esolu le paradoxe, puisque le Trilemme
CRO \'etablit qu'aucun protocole ne peut le r\'esoudre enti\`erement, que
ce soit en environnement classique~\cite{MINKA2026CRO} ou
quantique~\cite{MINKA2026PQCRO}. Il dispose toutefois d\'esormais de ce
que la cryptographie classique ne lui avait jamais offert, \`a savoir une
m\'etrique pour mesurer \`a quel point un protocole s'approche de ce qu'il
peut interpr\'eter~\cite{MINKA2026CLO} ainsi qu'une famille de constructions pour
commencer. Le paradoxe de l'authenticit\'e invisible n'est plus invisible
puisqu'il est formalis\'e et born\'e, sous des hypoth\`eses qu'il reste \`a mesurer. Ce que l'arbitre ne comprend pas
encore enti\`erement, la th\'eorie lui dit maintenant combien cela lui
co\^ute en bits.
\end{boiteArbitre}